%% ============================================================
%%  SECCIÓN 12 — Observaciones
%%  Responsable: Minilux
%% ============================================================

\section{Observaciones}
\label{sec:observaciones}

\begin{enumerate}

    \item \textbf{Punto de desprendimiento impreciso:} La adición de
    arena mediante pipeta no permite un control fino del instante exacto
    en que el anillo se desprende de la superficie, lo que introduce una
    sobreestimación sistemática de la fuerza de tensión superficial en
    cada medición.

    \item \textbf{Contaminación del agua:} El uso prolongado del mismo
    recipiente pudo introducir trazas de detergente u otras impurezas en
    el agua, reduciendo $\gamma$ respecto al valor de agua pura a
    \SI{20}{\celsius}.

    \item \textbf{Temperatura no controlada:} El coeficiente de tensión
    superficial del agua disminuye con la temperatura. No se registró la
    temperatura del agua durante el experimento, por lo que no es posible
    corregir el valor de referencia bibliográfico.

    \item \textbf{Evaporación y drenaje en el segundo método:} La
    película jabonosa comienza a evaporarse y drenar desde el momento en
    que se retira el dispositivo de la solución, alterando las dimensiones
    geométricas $a$, $b$ y $h$ antes de completar la medición.

    \item \textbf{Resolución discreta del brazo mayor:} Los jinetillos
    solo pueden colocarse en ranuras enteras, lo que limita la resolución
    de la masa efectiva y genera una incertidumbre irreducible en
    $\gamma_i$.

\end{enumerate}

%% ============================================================
%%  SECCIÓN 13 — Sugerencias
%%  Responsable: Minilux
%% ============================================================

\section{Sugerencias}
\label{sec:sugerencias}

\begin{enumerate}

    \item \textbf{Registrar la temperatura del agua:} Medir la
    temperatura del líquido antes y después de cada serie de mediciones
    permitiría corregir el valor de referencia
    $\gamma_{\text{ref}}(T)$ y reducir el error relativo.

    \item \textbf{Renovar el agua entre mediciones:} Usar agua destilada
    fresca en cada repetición evitaría la acumulación de impurezas y
    garantizaría condiciones reproducibles.

    \item \textbf{Emplear un tornillo micrométrico para la arena:}
    Sustituir la pipeta por un tornillo micrométrico o un sistema de
    goteo controlado permitiría detectar el punto de desprendimiento con
    mayor precisión.

    \item \textbf{Fotografiar la película jabonosa:} Capturar una imagen
    del dispositivo de tubitos inmediatamente después de retirarlo de la
    solución permitiría medir $a$, $b$ y $h$ con mayor precisión mediante
    análisis de imagen, eliminando el error por evaporación.

    \item \textbf{Repetir el ciclo al menos tres veces:} Realizar tres
    ciclos completos por cada método y promediar los resultados reduciría
    la incertidumbre estadística y permitiría identificar valores atípicos.

\end{enumerate}

%% ============================================================
%%  SECCIÓN 14 — Conclusiones
%%  Responsable: Minilux
%% ============================================================

\section{Conclusiones}
\label{sec:conclusiones}

\begin{enumerate}

    \item \textbf{Determinación de $\gamma$ por el primer método:} A
    partir del equilibrio de torques en la balanza de Mohr--Westphal se
    obtuvo un coeficiente de tensión superficial promedio de
    $\bar{\gamma}_1 = (39{,}73 \pm 0{,}88)\;\si{\milli\newton\per\meter}$,
    con un error relativo de \qty{45,4}{\percent} respecto al valor
    bibliográfico $\gamma_{\text{ref}} = \qty{72,8}{\milli\newton\per\meter}$
    \cite{serway}. La discrepancia se atribuye principalmente a la
    imprecisión en la detección del punto de desprendimiento del anillo
    y a la posible contaminación del agua.

    \item \textbf{Determinación de $\gamma$ por el segundo método:} La
    aplicación de la Ec.~\eqref{eq:gamma_m2} a los parámetros geométricos
    de la película jabonosa arrojó
    $\gamma_2 = \qty{12,21}{\milli\newton\per\meter}$, con un error
    relativo de \qty{83,2}{\percent}. Este valor, significativamente
    inferior al bibliográfico, es consistente con la reducción de tensión
    superficial producida por el surfactante presente en la solución
    jabonosa.

    \item \textbf{Comparación entre métodos:} El primer método, realizado
    con agua pura, produce un valor de $\gamma$ más próximo al
    referencial que el segundo. La diferencia entre ambos resultados
    constituye una verificación cualitativa del efecto de los
    surfactantes sobre la tensión superficial.

    \item \textbf{Sensibilidad geométrica del segundo método:} El
    análisis de la Ec.~\eqref{eq:gamma_m2} muestra que el resultado es
    altamente sensible a variaciones en $b$: un error de
    $\Delta b = \qty{0,1}{\milli\meter}$ se propaga como una variación
    relativa de $\sim\qty{6}{\percent}$ en $\gamma$, lo que hace de este
    método el menos robusto de los dos ante imprecisiones de medición.

    \item \textbf{Reproducibilidad del primer método:} La dispersión
    relativa de \qty{6}{\percent} entre las cinco mediciones válidas
    indica una reproducibilidad aceptable dentro de las limitaciones
    instrumentales de la balanza de Mohr--Westphal.

\end{enumerate}